\skripsifrontchapter{ABSTRAK}

\noindent\textbf{\namamahasiswa, \judulskripsi}\par
\noindent\textbf{Pembimbing: \pembimbingsatu{} dan \pembimbingdua}\par
\vspace{\baselineskip}

Naskah contoh ini memperlihatkan penggunaan template skripsi berbasis LaTeX untuk menata bagian awal, enam bab utama, daftar referensi, dan lampiran. Tujuannya adalah menyediakan titik awal yang dapat disunting serta membantu pengguna mengenali hubungan antara struktur sumber dan tampilan dokumen. Materi yang disertakan merupakan penjelasan format, bukan laporan penelitian empiris. Identitas mahasiswa dan pembimbing menggunakan penanda generik, sedangkan ruang persetujuan, pernyataan orisinalitas, tanggal, dan tanda tangan tetap memerlukan pengisian berdasarkan keadaan sebenarnya. Bab pendahuluan menunjukkan cara menghubungkan konteks, pertanyaan, tujuan, manfaat, dan batasan. Bab landasan memberikan contoh sitasi dokumentasi resmi serta definisi simbol matematika. Bab metodologi menampilkan hubungan tahap dan bahan pemeriksaan melalui tabel dan diagram vektor yang dibuat langsung di dalam naskah. Bab hasil memperlihatkan tabel ringkas, tabel yang melintasi halaman dengan judul kolom berulang, serta rujukan silang ke gambar dan persamaan. Isi tabel menggunakan nama komponen dan deskripsi keluaran tanpa skor, jumlah sampel, atau keberhasilan pengujian yang direka. Bab pembahasan menerangkan hubungan bukti, penafsiran, dan keterbatasan, sementara penutup menunjukkan bentuk simpulan serta saran yang perlu disesuaikan. Daftar referensi berisi sumber daring resmi untuk memperlihatkan pemanggilan bibliografi, dan lampiran menata catatan pendukung. Pengguna perlu mengganti seluruh materi contoh dengan isi penelitian yang dapat dipertanggungjawabkan, memeriksa aturan penulisan yang berlaku, serta membaca kembali keluaran setelah perubahan terakhir. Dokumen yang berhasil dikompilasi tetap memerlukan pemeriksaan isi, format, rujukan, dan kelengkapan administratif sebelum digunakan sebagai naskah akademik.

\noindent\textbf{Kata kunci:} LaTeX, template skripsi, struktur naskah, sitasi, rujukan silang, pemeriksaan format.
